% GENERATED FILE. Edit canonical lessons, then run npm run generate:books.
% canonical-source-hash: fnv1a64:da7bbbd0f08834cf
% canonical-lessons: KA-C496-hambalisu, KA-C496-halahalisu, KA-C496-bharavase, KA-C496-katara, KA-C496-nirala

\chapter{Hope and Longing}
\label{ch:ka-hope-and-longing}
\hlchaptermodality{496}

\begin{chapteropening}
\textbf{By the end of this chapter:} \emph{I can say to long for, to rue a loss, hope, eagerness and relieved.}

Complete the last lesson of chapter 496: I can say to long for, to rue a loss, hope, eagerness and relieved.
\end{chapteropening}

\section[\texorpdfstring{\kn{ಹಂಬಲಿಸು} (hambalisu)}{ಹಂಬಲಿಸು (hambalisu)}]{\kn{ಹಂಬಲಿಸು} (hambalisu) — to long for; to yearn}

\label{lesson:KA-C496-hambalisu}



\begin{quote}
Before the new one: say the Kannada for motherly love, then the Kannada for care.
\end{quote}

\subsection*{You'll want to know: \kn{ಹಂಬಲಿಸು}}
\textbf{\kn{ಹಂಬಲಿಸು}} — \emph{hambalisu} — \textquotedblleft{}to long for; to yearn\textquotedblright{}.

\begin{quote}
\kn{ವಿದೇಶದಲ್ಲಿದ್ದ} \kn{ಮಗ} \kn{ಅಮ್ಮನ} \kn{ಅಡುಗೆಗಾಗಿ} \kn{ಹಂಬಲಿಸುತ್ತಿದ್ದನು}.
\end{quote}

\begin{quote}
\emph{vidēśadallidda maga ammana aḍugegāgi hambalisuttiddanu.}
\end{quote}

\begin{quote}
The son living abroad used to long for his mother's cooking.
\end{quote}

Past (I …): \kn{ಹಂಬಲಿಸಿದೆ}, \emph{hambaliside}.

What you long for takes -\kn{ಗಾಗಿ}. The noun is \kn{ಹಂಬಲ}, a longing.

\begin{cousinweb}
\textbf{\kn{ಹಂಬಲಿಸು}} (\emph{hambalisu}) is native Dravidian: DEDR 3931 lists Kannada \emph{pambal}, \textquotedblleft{}ardent desire\textquotedblright{}, and \emph{pambalisu}, \textquotedblleft{}to desire ardently\textquotedblright{}. Old Kannada p- became h- in the modern language.
\end{cousinweb}

\begin{grammarlens}[title={What you've built}]
One more word for describing what you lived and felt.
\end{grammarlens}

\subsection*{Your turn}
\begin{itemize}
\raggedright
  \item \emph{Say it:} \emph{hambalisu}
  \item \emph{Say it:} \emph{hambalisu}, once more
  \item \emph{Say it:} \emph{vidēśadallidda maga ammana aḍugegāgi hambalisuttiddanu.} — The son living abroad used to long for his mother's cooking.
  \item \emph{Recall:} say the Kannada for motherly love, then the Kannada for care, then say \emph{hambalisu} again
\end{itemize}

\begin{tcolorbox}[breakable,colback=teal!4,colframe=teal!35!black,title={Before you move on}]
What does \kn{ಹಂಬಲಿಸು} mean? (\textquotedblleft{}To long for; to yearn\textquotedblright{}.) Say it once more.
\end{tcolorbox}
\section[\texorpdfstring{\kn{ಹಳಹಳಿಸು} (haḷahaḷisu)}{ಹಳಹಳಿಸು (haḷahaḷisu)}]{\kn{ಹಳಹಳಿಸು} (haḷahaḷisu) — to rue a loss; to pine for the past}

\label{lesson:KA-C496-halahalisu}



\begin{quote}
Before the new one: say the Kannada for care, then the Kannada for to long for.
\end{quote}

\subsection*{You'll want to know: \kn{ಹಳಹಳಿಸು}}
\textbf{\kn{ಹಳಹಳಿಸು}} — \emph{haḷahaḷisu} — \textquotedblleft{}to rue a loss; to pine for the past\textquotedblright{}.

\begin{quote}
\kn{ಮಾರಿದ} \kn{ಹಳೆಯ} \kn{ಮನೆಗಾಗಿ} \kn{ಅಜ್ಜ} \kn{ಹಳಹಳಿಸಿದರು}.
\end{quote}

\begin{quote}
\emph{mārida haḷeya manegāgi ajja haḷahaḷisidaru.}
\end{quote}

\begin{quote}
Grandfather pined for the old house that had been sold.
\end{quote}

Past (I …): \kn{ಹಳಹಳಿಸಿದೆ}, \emph{haḷahaḷiside}.

A longing mixed with regret; \kn{ಹಳಹಳಿಕೆ} is the feeling.

\begin{cousinweb}
\textbf{\kn{ಹಳಹಳಿಸು}} (\emph{haḷahaḷisu}) is a native Kannada verb for longing after something gone; its origin is not settled.
\end{cousinweb}

\begin{grammarlens}[title={What you've built}]
One more word for describing what you lived and felt.
\end{grammarlens}

\subsection*{Your turn}
\begin{itemize}
\raggedright
  \item \emph{Say it:} \emph{haḷahaḷisu}
  \item \emph{Say it:} \emph{haḷahaḷisu}, once more
  \item \emph{Say it:} \emph{mārida haḷeya manegāgi ajja haḷahaḷisidaru.} — Grandfather pined for the old house that had been sold.
  \item \emph{Recall:} say the Kannada for care, then the Kannada for to long for, then say \emph{haḷahaḷisu} again
\end{itemize}

\begin{tcolorbox}[breakable,colback=teal!4,colframe=teal!35!black,title={Before you move on}]
What does \kn{ಹಳಹಳಿಸು} mean? (\textquotedblleft{}To rue a loss; to pine for the past\textquotedblright{}.) Say it once more.
\end{tcolorbox}
\section[\texorpdfstring{\kn{ಭರವಸೆ} (bharavase)}{ಭರವಸೆ (bharavase)}]{\kn{ಭರವಸೆ} (bharavase) — hope; assurance}

\label{lesson:KA-C496-bharavase}



\begin{quote}
Before the new one: say the Kannada for to long for, then the Kannada for to rue a loss.
\end{quote}

\subsection*{You'll want to know: \kn{ಭರವಸೆ}}
\textbf{\kn{ಭರವಸೆ}} — \emph{bharavase} — \textquotedblleft{}hope; assurance\textquotedblright{}.

\begin{quote}
\kn{ವೈದ್ಯರ} \kn{ಮಾತಿನಿಂದ} \kn{ನಮಗೆ} \kn{ಭರವಸೆ} \kn{ಬಂದಿತು}.
\end{quote}

\begin{quote}
\emph{vaidyara mātininda namage bharavase banditu.}
\end{quote}

\begin{quote}
The doctor's words gave us hope.
\end{quote}

\kn{ಭರವಸೆ} \kn{ಕೊಡು}, to give an assurance, to promise.

\begin{cousinweb}
\textbf{\kn{ಭರವಸೆ}} (\emph{bharavase}) is usually traced to Hindustani \emph{bharosā}, \textquotedblleft{}trust, reliance\textquotedblright{}; the path into Kannada is not settled. In Kannada it is both hope and a promise you can lean on.
\end{cousinweb}

\begin{grammarlens}[title={What you've built}]
One more word for describing what you lived and felt.
\end{grammarlens}

\subsection*{Your turn}
\begin{itemize}
\raggedright
  \item \emph{Say it:} \emph{bharavase}
  \item \emph{Say it:} \emph{bharavase}, once more
  \item \emph{Say it:} \emph{vaidyara mātininda namage bharavase banditu.} — The doctor's words gave us hope.
  \item \emph{Recall:} say the Kannada for to long for, then the Kannada for to rue a loss, then say \emph{bharavase} again
\end{itemize}

\begin{tcolorbox}[breakable,colback=teal!4,colframe=teal!35!black,title={Before you move on}]
What does \kn{ಭರವಸೆ} mean? (\textquotedblleft{}Hope; assurance\textquotedblright{}.) Say it once more.
\end{tcolorbox}
\section[\texorpdfstring{\kn{ಕಾತರ} (kātara)}{ಕಾತರ (kātara)}]{\kn{ಕಾತರ} (kātara) — eagerness; anxious longing}

\label{lesson:KA-C496-katara}



\begin{quote}
Before the new one: say the Kannada for to rue a loss, then the Kannada for hope.
\end{quote}

\subsection*{You'll want to know: \kn{ಕಾತರ}}
\textbf{\kn{ಕಾತರ}} — \emph{kātara} — \textquotedblleft{}eagerness; anxious longing\textquotedblright{}.

\begin{quote}
\kn{ಫಲಿತಾಂಶಕ್ಕಾಗಿ} \kn{ವಿದ್ಯಾರ್ಥಿಗಳು} \kn{ಕಾತರದಿಂದ} \kn{ಕಾಯುತ್ತಿದ್ದರು}.
\end{quote}

\begin{quote}
\emph{phalitāṃśakkāgi vidyārthigaḷu kātaradinda kāyuttiddaru.}
\end{quote}

\begin{quote}
The students were waiting eagerly for the results.
\end{quote}

\kn{ಕಾತರದಿಂದ}, eagerly, impatiently.

\begin{cousinweb}
\textbf{\kn{ಕಾತರ}} (\emph{kātara}) is borrowed from Sanskrit \emph{kātara}, \textquotedblleft{}confused, discouraged, cowardly\textquotedblright{} (CDIAL 3027). In Kannada the nervous edge stayed and turned into impatient longing.
\end{cousinweb}

\begin{grammarlens}[title={What you've built}]
One more word for describing what you lived and felt.
\end{grammarlens}

\subsection*{Your turn}
\begin{itemize}
\raggedright
  \item \emph{Say it:} \emph{kātara}
  \item \emph{Say it:} \emph{kātara}, once more
  \item \emph{Say it:} \emph{phalitāṃśakkāgi vidyārthigaḷu kātaradinda kāyuttiddaru.} — The students were waiting eagerly for the results.
  \item \emph{Recall:} say the Kannada for to rue a loss, then the Kannada for hope, then say \emph{kātara} again
\end{itemize}

\begin{tcolorbox}[breakable,colback=teal!4,colframe=teal!35!black,title={Before you move on}]
What does \kn{ಕಾತರ} mean? (\textquotedblleft{}Eagerness; anxious longing\textquotedblright{}.) Say it once more.
\end{tcolorbox}
\section[\texorpdfstring{\kn{ನಿರಾಳ} (nirāḷa)}{ನಿರಾಳ (nirāḷa)}]{\kn{ನಿರಾಳ} (nirāḷa) — relieved; at ease}

\label{lesson:KA-C496-nirala}



\begin{quote}
Before the new one: say the Kannada for hope, then the Kannada for eagerness.
\end{quote}

\subsection*{You'll want to know: \kn{ನಿರಾಳ}}
\textbf{\kn{ನಿರಾಳ}} — \emph{nirāḷa} — \textquotedblleft{}relieved; at ease\textquotedblright{}.

\begin{quote}
\kn{ಪರೀಕ್ಷೆ} \kn{ಮುಗಿದ} \kn{ಮೇಲೆ} \kn{ನನಗೆ} \kn{ನಿರಾಳ} \kn{ಆಯಿತು}.
\end{quote}

\begin{quote}
\emph{parīkṣe mugida mēle nanage nirāḷa āyitu.}
\end{quote}

\begin{quote}
Once the exam was over, I felt relieved.
\end{quote}

\kn{ನಿರಾಳವಾಗಿ}, freely, with a light heart.

\begin{cousinweb}
\textbf{\kn{ನಿರಾಳ}} (\emph{nirāḷa}) is an everyday Kannada word for the ease that follows worry; its origin is not settled.
\end{cousinweb}

\begin{grammarlens}[title={What you've built}]
One more word for describing what you lived and felt.
\end{grammarlens}

\subsection*{Your turn}
\begin{itemize}
\raggedright
  \item \emph{Say it:} \emph{nirāḷa}
  \item \emph{Say it:} \emph{nirāḷa}, once more
  \item \emph{Say it:} \emph{parīkṣe mugida mēle nanage nirāḷa āyitu.} — Once the exam was over, I felt relieved.
  \item \emph{Recall:} say the Kannada for hope, then the Kannada for eagerness, then say \emph{nirāḷa} again
\end{itemize}

\begin{tcolorbox}[breakable,colback=teal!4,colframe=teal!35!black,title={Before you move on}]
What does \kn{ನಿರಾಳ} mean? (\textquotedblleft{}Relieved; at ease\textquotedblright{}.) Say it once more.
\end{tcolorbox}
